\paragraph{素数寄存器幂等锥、有限素窗口半圆维与二原子信息几何的闭合描述}
在素数寄存器估值算术化、链上规范投影的布尔化、有限素数支撑乘法单子的半圆维分离以及 bin-fold 相对均匀基线的普适 \(f\)-散度塌缩已经建立之后，还可进一步抽取出下列五条直接后果。

\begin{theorem}[素数寄存器幂等锥的秩分层、生成多项式与 Dirichlet 闭式]\label{thm:xi-prime-register-idempotent-cone-rank-dirichlet}
沿用定理 \ref{thm:killo-finite-semigroup-prime-valuation-arithmetization} 与定理 \ref{thm:xi-prime-register-idempotent-exact-count} 的记号，并记
\[
\mathcal S_n:=\left\{\prod_{i=0}^{n-1}q_i^{a_i}: a_i\in\{0,\dots,n-1\}\right\},
\]
\[
E_n:=\operatorname{Idem}(\mathcal S_n),
\qquad
E_{n,k}:=\operatorname{Idem}_k(\mathcal S_n),
\qquad
\operatorname{rk}(N):=|\operatorname{Im}(f_N)|.
\]
则对每个 \(1\le k\le n\) 都有
\[
|E_{n,k}|=\binom{n}{k}k^{n-k},
\qquad
|E_n|=\sum_{k=1}^{n}\binom{n}{k}k^{n-k}.
\]
其秩生成多项式满足
\[
\mathcal E_n(y):=\sum_{N\in E_n}y^{\operatorname{rk}(N)}
=
\sum_{k=1}^{n}\binom{n}{k}k^{n-k}y^k.
\]
此外，有限 Dirichlet 生成式具有显式闭式
\[
\mathfrak E_n(s):=\sum_{N\in E_n}N^{-s}
=
\sum_{\varnothing\neq F\subseteq\{0,\dots,n-1\}}
\left(\prod_{j\in F}q_j^{-sj}\right)
\prod_{i\notin F}\left(\sum_{j\in F}q_i^{-sj}\right).
\]
\end{theorem}
\begin{proof}
由推论 \ref{cor:killo-projection-idempotent-prime-register}，\(N\in E_n\) 当且仅当 \(f_N\) 为幂等映射，且此时
\[
\operatorname{Im}(f_N)=\operatorname{Fix}(f_N).
\]
因此，定理 \ref{thm:xi-prime-register-idempotent-exact-count} 已直接给出
\[
|E_{n,k}|=\binom{n}{k}k^{n-k},
\qquad
|E_n|=\sum_{k=1}^{n}\binom{n}{k}k^{n-k}.
\]
将秩 \(k\) 的计数按 \(y^k\) 加权求和，即得 \(\mathcal E_n(y)\) 的闭式。

再求 Dirichlet 生成式。固定非空子集 \(F\subseteq\{0,\dots,n-1\}\)。由
\[
\operatorname{Im}(f_N)=\operatorname{Fix}(f_N)=F
\]
可知，此类幂等元恰由下列独立条件刻画：
\[
v_{q_j}(N)=j\qquad (j\in F),
\qquad
v_{q_i}(N)\in F\qquad (i\notin F).
\]
因此，对 \(j\in F\) 必有贡献因子 \(q_j^{-sj}\)；而对 \(i\notin F\)，其指数可在 \(F\) 中任取，故贡献
\[
\sum_{j\in F}q_i^{-sj}.
\]
各坐标彼此独立，相乘后即得到固定像集 \(F\) 的全部贡献
\[
\left(\prod_{j\in F}q_j^{-sj}\right)
\prod_{i\notin F}\left(\sum_{j\in F}q_i^{-sj}\right).
\]
再对全部非空 \(F\) 求和，即得所示 Dirichlet 闭式。证毕。
\end{proof}

\begin{theorem}[单调幂等骨架的合成退化与指数稀疏性]\label{thm:xi-monotone-idempotent-skeleton-gcd-lcm-sparsity}
沿用定理 \ref{thm:killo-chain-interior-godel-gcd-lcm} 的记号，记
\[
\mathcal I_n:=\mathrm{Int}(C_n),
\qquad
\mathcal I_{n,k}:=\{I\in\mathcal I_n:\ |\mathrm{Fix}(I)|=k\}
\qquad(1\le k\le n).
\]
则对任意满足 \(0\in F\cap G\subseteq\{0,\dots,n-1\}\) 的子集 \(F,G\)，有
\[
I_F\circ I_G=I_{F\cap G}=I_G\circ I_F,
\]
从而
\[
\kappa(I_F\circ I_G)=\gcd(\kappa(I_F),\kappa(I_G)),
\qquad
\kappa(I_F\vee I_G)=\mathrm{lcm}(\kappa(I_F),\kappa(I_G)).
\]
特别地，\(\mathcal I_n\) 在函数合成下构成交换幂等半群，并经 \(\kappa\) 与平方自由整数上的 \(\gcd\)-半群同构。

此外，
\[
|\mathcal I_{n,k}|=\binom{n-1}{k-1},
\]
并且其在全部秩 \(k\) 幂等元中所占比例满足
\[
\frac{|\mathcal I_{n,k}|}{|E_{n,k}|}
=
\frac{\binom{n-1}{k-1}}{\binom{n}{k}k^{n-k}}
=
\frac{k}{n}\,k^{-(n-k)}.
\]
因而只要 \(1\le k\le n-1\)，该比例便随 \(n-k\) 指数衰减。
\end{theorem}
\begin{proof}
由定理 \ref{thm:killo-chain-interior-godel-gcd-lcm}，任意 \(I\in\mathcal I_n\) 都唯一写成
\[
I=I_F,
\qquad
F=\mathrm{Fix}(I),
\qquad
0\in F.
\]
对任意 \(i\in\{0,\dots,n-1\}\)，直接计算得
\[
(I_F\circ I_G)(i)
=
I_F\!\bigl(\max(G\cap\{0,\dots,i\})\bigr)
=
\max\bigl(F\cap G\cap\{0,\dots,i\}\bigr)
=
I_{F\cap G}(i).
\]
由于交运算交换，故
\[
I_F\circ I_G=I_G\circ I_F=I_{F\cap G},
\]
从而 \(\mathcal I_n\) 在合成下为交换幂等半群。再由同一定理的算术化公式，
\[
\kappa(I_{F\cap G})=\gcd(\kappa(I_F),\kappa(I_G)),
\qquad
\kappa(I_F\vee I_G)=\mathrm{lcm}(\kappa(I_F),\kappa(I_G)),
\]
即得第一部分。

对计数部分，只需从 \(\{1,\dots,n-1\}\) 中选出 \(k-1\) 个元素并与强制出现的 \(0\) 合并为固定点集 \(F\)，故
\[
|\mathcal I_{n,k}|=\binom{n-1}{k-1}.
\]
再与定理 \ref{thm:xi-prime-register-idempotent-cone-rank-dirichlet} 中
\[
|E_{n,k}|=\binom{n}{k}k^{n-k}
\]
相除，便得到所示比例公式。最后一条由 \(k\ge 1\) 时 \(k^{-(n-k)}\) 的指数衰减立即成立。证毕。
\end{proof}

\begin{theorem}[有限素窗口同态半圆维的严格线性差距]\label{thm:xi-finite-prime-window-halfcircle-linear-gap}
取互异素数 \(p_1,\dots,p_r\)，并记
\[
M_r:=\langle p_1,\dots,p_r\rangle
=
\left\{p_1^{a_1}\cdots p_r^{a_r}: a_1,\dots,a_r\in\ZZ_{\ge 0}\right\}.
\]
则
\[
\cdim^{\mathrm{hom}}_{+}(M_r)=\frac{r}{2},
\qquad
\cdim^{\mathrm{impl}}(M_r)=\frac12.
\]
特别地，
\[
\cdim^{\mathrm{hom}}_{+}(M_r)-\cdim^{\mathrm{impl}}(M_r)=\frac{r-1}{2},
\]
故有限素窗口上的结构层与实现层半圆维差距已经呈严格线性增长。
\end{theorem}
\begin{proof}
令
\[
S_r:=\{p_1,\dots,p_r\}\subset\mathbb P.
\]
则
\[
M_r=\NN_{S_r}^\times.
\]
所述等式正是定理 \ref{thm:xi-finite-prime-support-multiplicative-half-circle-dimension} 在 \(S=S_r\) 时的直接转写。证毕。
\end{proof}

\begin{theorem}[单一 \texorpdfstring{$\chi^2$}{chi2} 极限常数完备编码全部 Csisz\'ar 几何]\label{thm:xi-chi2-complete-csiszar-geometry}
沿用命题 \ref{prop:fold-bin-f-divergence-collapse} 的记号，并记
\[
f_2(t):=(t-1)^2,
\qquad
\mathcal C_2:=D_{f_2}^\infty=\lim_{m\to\infty}\chi^2(\mu_m\Vert u_m).
\]
则
\[
\varphi=\bigl(5(\mathcal C_2+1)-1\bigr)^{1/3}.
\]
若进一步置
\[
\xi:=\bigl(5(\mathcal C_2+1)-1\bigr)^{1/3},
\]
则对任意凸函数 \(f:(0,\infty)\to\RR\) 满足 \(f(1)=0\)，其 bin-fold 极限 \(f\)-散度都满足
\[
D_f^\infty
=
\frac{1}{\xi}\,f\!\left(\frac{\xi^2}{\sqrt5}\right)
+\frac{1}{\xi^2}\,f\!\left(\frac{\xi}{\sqrt5}\right).
\]
因此，bin-fold 相对均匀基线的全部渐近 Csisz\'ar 信息几何并非无限维数据，而由单一标量 \(\mathcal C_2\) 完整决定。
\end{theorem}
\begin{proof}
由命题 \ref{prop:fold-bin-f-divergence-collapse} 取 \(f=f_2\)，得到
\[
\mathcal C_2
=
\frac{1}{\varphi}\left(\frac{\varphi^2}{\sqrt5}\right)^2
+\frac{1}{\varphi^2}\left(\frac{\varphi}{\sqrt5}\right)^2
-1
=
\frac{\varphi^3+1}{5}-1.
\]
整理即得
\[
\varphi^3=5(\mathcal C_2+1)-1,
\]
从而
\[
\varphi=\bigl(5(\mathcal C_2+1)-1\bigr)^{1/3}.
\]
将
\[
\xi=\bigl(5(\mathcal C_2+1)-1\bigr)^{1/3}=\varphi
\]
代回命题 \ref{prop:fold-bin-f-divergence-collapse} 的两点闭式，便得到任意 \(f\) 的极限表达式。故全部 \(D_f^\infty\) 都由 \(\mathcal C_2\) 唯一决定。证毕。
\end{proof}

\begin{theorem}[渐近 \texorpdfstring{$f$}{f}-散度泛函的二原子刚性]\label{thm:xi-asymptotic-f-divergence-functional-two-atom-rigidity}
记
\[
\nu_\varphi
:=
\frac{1}{\varphi}\,\delta_{\varphi^2/\sqrt5}
+\frac{1}{\varphi^2}\,\delta_{\varphi/\sqrt5}.
\]
设 \(K\subset(0,\infty)\) 为包含
\(
\{\varphi/\sqrt5,\varphi^2/\sqrt5\}
\)
的紧区间，\(\nu\) 为 \(K\) 上的 Borel 概率测度。若对一切 \(\alpha>1\) 都有
\[
\int_K \bigl(t^\alpha-1-\alpha(t-1)\bigr)\,\dd\nu(t)
=
D_{f_\alpha}^\infty
=
\int_K \bigl(t^\alpha-1-\alpha(t-1)\bigr)\,\dd\nu_\varphi(t),
\]
其中
\[
f_\alpha(t):=t^\alpha-1-\alpha(t-1),
\]
则必有
\[
\nu=\nu_\varphi.
\]
换言之，bin-fold 的渐近 \(f\)-散度泛函在测度论层面唯一对应于该二原子模型。
\end{theorem}
\begin{proof}
设
\[
\sigma:=\nu-\nu_\varphi.
\]
对 \(z\in\CC\) 定义
\[
G(z):=\int_K \bigl(t^z-1-z(t-1)\bigr)\,\dd\sigma(t),
\qquad
t^z:=e^{z\log t}.
\]
由于 \(K\subset(0,\infty)\) 为紧集，\(\log t\) 在 \(K\) 上有界，故 \(G\) 是整函数。由假设，对每个实数 \(\alpha>1\) 都有
\[
G(\alpha)=0.
\]
零点集含有聚点，因此由恒等定理可知
\[
G\equiv 0 \qquad \text{于 } \CC.
\]
于是对一切 \(z\in\CC\) 都有
\[
0=G''(z)=\int_K (\log t)^2 t^z\,\dd\sigma(t).
\]
定义有限符号测度
\[
\dd\lambda(t):=(\log t)^2 t^2\,\dd\sigma(t).
\]
则对任意整数 \(n\ge 0\)，
\[
\int_K t^n\,\dd\lambda(t)=G''(n+2)=0.
\]
因此 \(\lambda\) 对 \(K\) 上所有多项式的积分皆为零。由 Stone--Weierstrass 定理，多项式在 \(C(K)\) 中稠密，故
\[
\lambda=0.
\]
而对 \(t\neq 1\) 有
\[
(\log t)^2 t^2>0,
\]
故 \(\sigma\) 只能支撑在 \(\{1\}\) 上。另一方面，\(\nu\) 与 \(\nu_\varphi\) 都是概率测度，故
\[
\sigma(K)=0.
\]
因此 \(\sigma\) 不可能在 \(\{1\}\) 上留下任何非零质量，只能有
\[
\sigma=0.
\]
于是 \(\nu=\nu_\varphi\)。证毕。
\end{proof}

\paragraph{素数寄存器骨架与二原子信息几何的双侧闭合}
上述五条结果表明，素数寄存器侧的幂等结构已经可由秩分层、Dirichlet 级数与布尔骨架的指数稀疏性完全控制；与此同时，bin-fold 相对均匀基线的渐近信息几何则最终塌缩为由单一 \(\chi^2\) 常数编码的唯一二原子测度。两侧共同给出一条算术外置与信息塌缩相互对照的闭合描述。

\endinput
